\scanpage{11}
\mspage{5}
% archive foliation: 10
\diagram{\left\{
\begin{array}{llll}
&\text{l'image inverse de}&\mathbb{F}_2\times\{0\}&\text{est }V_S\\
&\rule{4cm}{0.3pt}&\{0\}\times\mathbb{F}_2&\text{est }V_A\\
\text{donc}&\rule{4cm}{0.3pt}&
\begin{array}{c}\text{diagonale}\\\text{de }\mathbb{F}_2\times\mathbb{F}_2\end{array}
&\text{est }\mathfrak{D}^+
\end{array}\right.}
ou encore \deleted{\unsure{ainsi}}
\diagram{\left\{
\begin{array}{lll}
\text{l'image inverse de}&(1,0)&\text{est }C\subset\unsure{\mathfrak{D}^*|\mathfrak{D}}\\
\rule{4cm}{0.3pt}&(0,1)&\text{est }D\subset\unsure{\mathfrak{D}^*|\mathfrak{D}}\\
\rule{4cm}{0.3pt}&(1,1)&\text{est }\underline{\omega}_Q\subset\unsure{\mathfrak{D}^*|\mathfrak{D}}
\end{array}\right.}

Les sous-groupes \unsure{d'ordre deux} de $\mathfrak{D}$ correspondent aux
éléments d'ordre 2, il y a donc les éléments

a) l'antipodisme $\underline{a}$, engendrant le centre $\{1,\underline{a}\}$

b) Les deux éléments de $C=V_S\setminus\{1,a\}$, engendrant resp. les deux
droites-coordonnées de $V_S\simeq\mathbb{F}_2^C$

c) Les deux éléments de $D=V_A\setminus\{1,a\}$, engendrant resp. les deux
droites coordonnées de $V_A\simeq\mathbb{F}_2^D$.

\margin{NB. ces quatre éléments n'étant pas centraux, les sous-groupes d'ordre 2
correspondants ne sont pas distingués ; leurs normalisateurs sont resp.
$V_S$ et $V_A$.
Donc parmi les huit ss-groupes de $\mathfrak{D}$ $\ne\{1\},\mathfrak{D}$,
il y en a 4 qui sont distingués
($V_S,V_A,\mathfrak{D}^+,\{1,\underline{a}\}$) et quatre qui ne le sont pas.}

(NB les \deleted{\unsure{autres}} 3 éléments de
$\unsure{\mathfrak{D}^*|\mathfrak{D}}$ sont d'ordre 1, et d'ordre 4 :
\[
\mathfrak{D}\simeq
\underbrace{\{1\}\amalg\{a\}\amalg C\amalg D\amalg\underline{\omega}_Q}
_{\text{déc. de }\mathfrak{D}\text{ en orbites sous }\operatorname{int}(\mathfrak{D})}
\qquad)
\]

On a donc fait la liste complète des sous-groupes de $\mathfrak{D}$ différents
de $\{1\}$ et de $\mathfrak{D}$ lui-$=$, il y en a trois d'ordre 4 i.e.
d'indice 2, savoir $V_S,V_A$ et $\mathfrak{D}^+$,
et cinq d'ordre 2 savoir le centre $\{1,\underline{a}\}$,
et les 4 sous-groupes \inserted{2 paquets de 2} d'ordre 2 distincts de ce dernier,
contenus respectivement dans $V_S$ et dans $V_A$.

On a associé à un carré $Q=(S,A,R)$ trois ensembles remarquables de cardinal 2,
savoir
\[
\unsure{B|D},\ \unsure{S|C},\ \underline{\omega}_Q
\]
\margin{\unsure{4}}
